\section{Pontryagin duality}

\begin{definition}[Canonical evaluation map]
\label{def:pd-eval}
\lean{PontryaginDuality.IsEvaluationIso,PontryaginDuality.evalAt,PontryaginDuality.evalCMH}
\leanok
The canonical evaluation homomorphism sends $a\in A$ to the character $\chi\mapsto\chi(a)$ on the Pontryagin dual.
\end{definition}

\begin{theorem}[Evaluation is bijective]
\label{thm:pd-bijective}
\lean{PontryaginDuality.evalCMH_bijective}
\notready
\uses{def:pd-eval}
For a locally compact Hausdorff abelian group, the evaluation map to the double dual is bijective. The required separation and Fourier-inversion ingredients are still a \texttt{sorry}.
\end{theorem}

\begin{theorem}[Evaluation is open]
\label{thm:pd-open}
\lean{PontryaginDuality.evalCMH_isOpenMap}
\notready
\uses{def:pd-eval}
The evaluation map is open; this step is also still a \texttt{sorry}.
\end{theorem}

\begin{theorem}[Pontryagin duality]
\label{thm:pontryagin-duality}
\lean{PontryaginDuality.MainTheorem}
\notready
\uses{thm:pd-bijective,thm:pd-open}
Every locally compact Hausdorff abelian group is canonically topologically isomorphic to its double Pontryagin dual.
\end{theorem}
